\documentclass[11pt]{article}
\usepackage[margin=1in]{geometry}
\usepackage[T1]{fontenc}
\usepackage[utf8]{inputenc}
\usepackage{lmodern}
\usepackage{hyperref}

\begin{document}

\section*{Manuscript Review}

\subsection*{Overall assessment}
This is a strong proposal draft with a clear research motivation, a coherent pipeline, and a useful appendix linking methods to implementation. The most important issues are methodological calibration (what the current design can and cannot claim), one implementation-level plotting error in the appendix code, and several places where claim strength should be reduced to match currently available evidence.

Technical coverage checks were completed for symbol definitions, sign/branch conventions, dimensional consistency, and implementation ambiguity. The manuscript does not use inverse-trigonometric or coordinate-transform formulas, so branch-convention risk is not central here; no concrete sign/dimensional equation errors were identified in the statistical expressions that are present (\(\chi^2\), \(\alpha\)).

\subsection*{Most important technical findings}

\begin{itemize}

\item \textbf{Unsupported-claim risk: hypothesis adjudication is stated before the required NRC analysis is complete.}\\
\textbf{Location:} Section ``Expected Outcomes and Alternative Hypotheses'' (statement that preliminary \(\chi^2\) analysis will determine which hypothesis is supported), read together with Section ``Conclusion'' and Appendix note that NRC semantic tagging is still pending.\\
\textbf{Problem:} The manuscript states that preliminary analysis will determine the supported hypothesis, but the hypotheses are defined around NRC--SCC comparison and NRC semantic labeling has not yet been run.\\
\textbf{Why it matters:} This overstates current evidential status and blurs proposal-stage versus result-stage conclusions.\\
\textbf{Suggested fix:} Rephrase to planned analysis language, e.g., ``After NRC semantic tagging, the preregistered \(\chi^2\)/exact-test pipeline will be used to evaluate H1--H3.''

\item \textbf{Likely internal-validity issue: register is confounded with grammatical mood.}\\
\textbf{Location:} Introduction (research question framed as register discrimination), Section ``Synthetic Control Corpus (SCC)'' condition definitions, and Section ``Expected Outcomes'' interpretation scope.\\
\textbf{Problem:} Conditions bind register and mood (HIGH=questions, NEUTRAL=imperatives, LOW=statements), so effects cannot be attributed cleanly to register alone. This limitation is acknowledged later, but downstream hypothesis wording still reads as primarily register-causal.\\
\textbf{Why it matters:} It weakens causal interpretability of the primary claim (register-specific idiomaticity differences).\\
\textbf{Suggested fix:} Either (a) reframe all hypotheses as ``condition effects'' rather than ``register effects'', or (b) use a crossed design (register \(\times\) mood) and model both factors.

\item \textbf{Method--implementation mismatch around sparse-cell handling.}\\
\textbf{Location:} Section ``Statistical Analysis'' (Fisher's Exact Test for small expected frequencies) versus Appendix code listing \texttt{visualizer.py}.\\
\textbf{Problem:} The text promises an exact-test fallback for sparse counts, but the provided analysis script only runs \texttt{chi2\_contingency} and no exact procedure.\\
\textbf{Why it matters:} Reproducibility suffers: readers cannot recover the stated inferential protocol from the documented code path.\\
\textbf{Suggested fix:} Specify and implement one concrete sparse-cell strategy (e.g., Monte-Carlo exact p-value for \(r\times c\) tables, or preregistered pairwise 2\(\times\)2 Fisher tests with multiplicity correction), then align manuscript and appendix.

\item \textbf{Definite implementation error in figure labeling (appendix script).}\\
\textbf{Location:} Appendix Section ``Python Code Documentation'', \texttt{visualizer.py}, line with \texttt{ax.bar\_label(c, fmt='\%.0f\%\%')}.\\
\textbf{Problem:} Bars are plotted from proportions in \([0,1]\), but labels are formatted as if values were percentages; e.g., 0.98 is displayed as ``1\%'' instead of ``98\%''.\\
\textbf{Why it matters:} This can materially misreport effect size in the plotted output and undermine interpretation of the appendix figures.\\
\textbf{Suggested fix:} Multiply by 100 before labeling or use explicit custom labels, e.g., \texttt{labels=[f"\{100*v:.0f\}\%" for v in values]}.

\item \textbf{Notation/terminology drift risk (NRC vs. TEC role labeling).}\\
\textbf{Location:} Introduction and Methodology subsection ``Natural Reference Corpus (NRC)'', plus the corpus representativeness discussion.\\
\textbf{Problem:} ``NRC'' (analysis dataset) and ``TEC'' (source corpus) are mostly used correctly, but some sentences switch baseline wording in ways that can blur whether claims refer to the sampled NRC subset or the broader TEC resource.\\
\textbf{Why it matters:} This creates implementation-facing ambiguity for replication (what exactly is compared in each test).\\
\textbf{Suggested fix:} Add one explicit definitional sentence early in Methodology: ``TEC is the source corpus; NRC denotes the extracted B1/B2 verb-instance dataset used in all statistical tests,'' then consistently reference NRC for inferential claims.

\end{itemize}

\subsection*{Meaningful editorial findings}

\begin{itemize}

\item \textbf{Over-strong wording in conclusion relative to measured outcomes.}\\
\textbf{Location:} Conclusion paragraph discussing learnability and processing load benefits.\\
\textbf{Problem:} The current wording reads close to an empirical learning-effect conclusion, while the study currently measures distributional usage categories, not acquisition or cognitive load outcomes.\\
\textbf{Why it matters:} It risks interpretive overreach in a proposal-stage document.\\
\textbf{Suggested fix:} Hedge to mechanism-compatible language.\\
\textbf{Candidate rewrite:} ``These distributional differences may indicate input properties associated with learnability, but direct effects on acquisition and cognitive load require separate learner-based experiments.''

\item \textbf{Reference-list hygiene and metadata consistency.}\\
\textbf{Location:} \texttt{bibliography.bib} (multiple entries).\\
\textbf{Problem:} Inconsistent capitalization in publisher names (e.g., ``Oxford university press''), plus local \texttt{file=} paths from a personal machine in bibliography entries.\\
\textbf{Why it matters:} Reduces professionalism and can leak irrelevant local metadata into shared academic artifacts.\\
\textbf{Suggested fix:} Normalize proper-name capitalization and remove non-portable local file-path fields before submission.

\item \textbf{Several high-impact sentence-level grammar defects in method/results framing.}\\
\textbf{Location:} Statistical analysis and hypothesis sections.\\
\textbf{Problem:} Phrases such as ``will additionally be used as to check'' and ``will show pattern similar to'' are ungrammatical at key technical points.\\
\textbf{Why it matters:} These occur in inferential-method statements, where clarity is essential for replicability.\\
\textbf{Suggested fix:} Tighten these lines for grammatical precision.\\
\textbf{Candidate rewrites:}\\
(a) ``Therefore, an exact-test procedure will be used when expected cell counts are small.''\\
(b) ``H2 predicts that the NRC will show \emph{a} pattern similar to SCC-NEUTRAL.''

\end{itemize}

\subsection*{Proofing-sweep items (from pattern scan and spot-check)}

\begin{itemize}
\item \textbf{Location:} SCC condition labels in the numbered list.\\
\textbf{Problem:} Semicolon-capitalization pattern inside parentheticals (e.g., ``Workplace; Questions'') reads as punctuation noise rather than clause separation.\\
\textbf{Suggested fix:} Replace semicolons with commas or an em dash for cleaner category notation.

\item \textbf{Location:} Idiom Principle example sentence.\\
\textbf{Problem:} Mid-sentence capitalization inconsistency in ``I got on the Bus.''\\
\textbf{Suggested fix:} Use ``I got on the bus.''

\item \textbf{Location:} Corpus representativeness paragraph.\\
\textbf{Problem:} Comma splice-like interruption in ``The NRC derived from an established pedagogical corpus (TEC), represents ...''\\
\textbf{Suggested fix:} Remove the comma after ``(TEC)''.
\end{itemize}

\subsection*{Highest-priority fixes}

\begin{itemize}
\item Align claim strength with analysis stage: do not state that hypotheses are already adjudicated before NRC semantic tagging is complete.
\item Resolve the register--mood confound in either design (crossed factors) or interpretation language (condition effects only).
\item Fix the percentage-label bug in \texttt{visualizer.py} before presenting appendix figures as quantitative evidence.
\item Harmonize sparse-cell testing plan between manuscript text and executable analysis code.
\end{itemize}

\subsection*{Items that need external verification}

\begin{itemize}
\item Comparative model-version examples in the limitations section (e.g., specific product-version names) should be verified against stable, citable release documentation or generalized to model classes.
\item The claim that B1/B2 represents the majority of adult ESL learners should be tied to a concrete source (policy report, enrollment statistics, or large-scale survey).
\item Opening broad claims about public uptake of specific chat systems should be supported with direct, recent adoption sources rather than indirect sector-level AI references.
\end{itemize}

\end{document}